\documentclass{article}
\usepackage{booktabs}
\usepackage{tabularx}
\usepackage[margin=1in]{geometry}
\begin{document}

\begin{table}[tbp] \centering
\newcolumntype{R}{>{\raggedleft\arraybackslash}X}
\newcolumntype{L}{>{\raggedright\arraybackslash}X}
\newcolumntype{C}{>{\centering\arraybackslash}X}

\caption{National crop production: structural-break period means, 1837-1915 (Phase C.14)}
\label{tab:quantity_periods}
\begin{tabularx}{\linewidth}{@{}lCCCCC@{}}

\toprule
&{(1)}&{(2)}&{(3)}&{(4)}&{(5)} \tabularnewline \midrule
{Period}&{Years (N)}&{Wine (1000 hl)}&{Potato (1000 q)}&{Cereal (1000 q)}&{Fruit (1000 q)} \tabularnewline
\midrule \addlinespace[\belowrulesep]
Pre-phylloxera (1837--1862)&26&1418&5849&4301&2225
Phylloxera era (1875--1895)&20&1363&7978&3477&3537
Recovery + ban (1895--1915)&21&1014&6849&2673&3991
\bottomrule \addlinespace[\belowrulesep]

\end{tabularx}
\\ \parbox{\linewidth}{\footnotesize Notes: Phase C.14 (verify\\_reconstruct\\_expand handoff). National crop production volumes from HSSO I.21a (1837-1991, yearly), summarized by structural-break period. Pillar 2 quantity-side corroboration of the substrate-substitution framework: framework predicts wine production DOWN and substitute production (potato, cereal) UP during the phylloxera era. Wine units are 1000 hl; crop units are 1000 quintals (q). Pre-phylloxera baseline starts 1837 (first data year in I.21a). Source: HSSO I.21a citing Mottu-Weber 1989 + Swiss Farmers' Secretariat statistics.}
\end{table}
\end{document}
